\clearpage\fancyhead[L]{\small JUNO.}\hypertarget{native-330-lower}{}\label{native-330-lower}
\begin{center}\textbf{JUNO.}\par\vspace{4pt}\hrule\vspace{4pt}\fontsize{10}{12}\selectfont Göttingische gelehrte Anzeigen. Issues 92. 93. pp. 913\,..\,914. 10 June 1811.\par\vspace{4pt}\hrule\vspace{4pt}\end{center}

\par\indent The last news about Juno can be found in issue 136 of our Notices for 1808 [p. 305 of this volume]. It gives Professor \textsc{Gauss}’s observations of this planet at our observatory in 1808, the determination of its fourth opposition, and the eighth set of elements. The fifth opposition, which occurred in January 1810, was observed nowhere. This is the first time that the determination of an opposition of one of the four new planets has been lost through the collective neglect of all astronomers. It was all the more important to ensure that the \SourceLeaf{331}\emph{sixth} opposition, which fell in April of this year, did not also go unobserved. The few observations Professor \textsc{Gauss} made, using only the ring micrometer because the planet was so extremely faint, will therefore be all the more valuable, since by all appearances they will be the only observations made anywhere on this occasion. For comparison, $\varphi$ Librae and several other nearby stars from the Histoire céleste were used. The results below will still need a small correction once the positions of these stars are determined more precisely.
\Needspace{10\baselineskip}
\begin{center}\fontsize{10}{12}\selectfont\setlength{\tabcolsep}{4pt}\renewcommand{\arraystretch}{1.15}\begin{tabular}{c|c|c|c}
1811 & \shortstack{mean time\\in Göttingen} & \shortstack{Right\\ascension} & \shortstack{Southern\\declination}\\\hline
22 April & $9^{\text{u}}\;51^{\text{m}}\;35^{\text{s}}$ & $216^\circ\;41'\;50''1$ & $0^\circ\;58'\;16''$\\
24 & $10\;32\;55$ & $216\;17\;58.7$ & $0\;46\;50$\\
25 & $10\;19\;22$ & $216\;6\;3.6$ & $0\;40\;4$\\
26 & $10\;17\;30$ & $215\;54\;7.3$ & $0\;33\;51$\\
\end{tabular}\end{center}

\par\indent The declination on 22 April is doubtful, and the declination on 24 April is also less reliable than the next two, when the planet was in a more favorable position.

\Needspace{8\baselineskip}\par\indent From these observations Professor \textsc{Gauss} obtained the following result for the opposition.
\begin{center}1811, 24 April. $19^{\text{u}}\;20^{\text{m}}\;12^{\text{s}}$ mean time in Göttingen\\$214^\circ\;8'\;48''3$ true heliocentric longitude\\$12\;55\;2.0$ true geocentric latitude, north\end{center}
\Needspace{20\baselineskip}\par\indent The large discrepancy between this position and the elements based on the first four oppositions now reveals, for Juno too, the effect of the perturbations caused especially by Jupiter. Calculating them for Juno will be as enormous a task as for Pallas. It is only right that this work should first be finished for Pallas, which was discovered earlier. Professor \textsc{Gauss} has therefore been content for the time being to base new elliptic elements on the oppositions of 1806, 1807, 1808, 1811. We give these elements below.

\begin{center}\fontsize{10}{12}\selectfont\setlength{\tabcolsep}{4pt}\renewcommand{\arraystretch}{1.15}\begin{tabular}{llr}
\shortstack[l]{Epoch of mean longitude, 1811,\\Göttingen meridian} & .\;. & $177^\circ\;48'\;1''8$\\
Mean daily tropical motion & .\;.\;.\;.\;.\;.\;.\;.\;.\;. & $813''2486$\\
Longitude of perihelion 1811 & .\;.\;.\;.\;.\;.\;.\;.\;.\;.\;.\;.\;.\;. & $53^\circ\;14'\;32''4$\\
Longitude of the ascending node 1811 & .\;.\;.\;.\;.\;.\;.\;.\;.\;. & $171\;9\;13.5$\\
Inclination of the orbit & .\;.\;.\;.\;.\;.\;.\;.\;.\;.\;.\;.\;.\;.\;.\;.\;. & $13\;4\;27.0$\\
Eccentricity & .\;.\;.\;.\;.\;.\;.\;.\;.\;.\;.\;.\;.\;.\;.\;.\;. & $=\sin14\;44\;9.1$\\
Logarithm of the semimajor axis & .\;.\;.\;.\;.\;.\;.\;.\;.\;.\;. & $0.4265711.$\\
\end{tabular}\end{center}
